\section{治理与法规对接}

\subsection{政策与合规需求}

GenAI/Agent 在企业中经常触及金融、医疗、政府等敏感领域。需要同时满足以下监管要求：

\begin{itemize}
    \item \textbf{数据主权}：数据必须在指定区域处理，不能跨境泄露
    \item \textbf{透明性}：用户要能看到为什么 agent 做出某个推荐
    \item \textbf{安全审计}：所有 tool 调用需记录并可供合规团队审查
    \item \textbf{偏见检测}：agent 不能在种族/性别上产生歧视性输出
\end{itemize}

\subsection{技术与治理结合}

合规不是单靠文档，而是要把 guardrail 编成 code：

\begin{itemize}
    \item policy-as-code：把规则写成可执行 policy（如 Open Policy Agent）
    \item simulation sandbox：每次改 prompt 先跑 sandbox 测试，确保不会违反 policy
    \item kill-switch：当检测到攻击/滥用时，agent 可自动进入安全 mode
\end{itemize}

\begin{warningbox}{监管不是 checkbox，而是持续对话}
监管机构会根据实际应用不断提出新要求。企业需要把 regulator view 嵌入 daily standup—当出现 new capability 时，先评估是否触发新的 policy，再上线。
\end{warningbox}

\subsection{本章小结}

治理需要跨组织协作：法律、合规、AI 安全、工程共同构建 policy-as-code、杀手开关和可解释的审计 trail。

